\section{Where to go from here}
\label{sec:practice}

\subsection{Software and a first experiment}

Everything in this tutorial can be reproduced on a laptop. The standard open-source tools are the following.

\begin{itemize}
\item \textbf{TransformerLens} \citep{nanda2022transformerlens} loads open models and gives access to every
intermediate activation through named \emph{hooks}. It is the usual choice for small models.
\item \textbf{nnsight} \citep{fiottokaufman2025nnsight} offers the same kind of access for arbitrary PyTorch
models, including large ones that run on a remote server.
\item \textbf{SAELens} \citep{bloom2024saelens} trains sparse autoencoders and loads published ones, such as
Gemma Scope \citep{lieberum2024gemma}.
\item \textbf{Neuronpedia} \citep{lin2023neuronpedia} is a website on which the latents of published
dictionaries can be browsed, tested on one's own text and used for steering. It also hosts attribution graphs.
\end{itemize}

The listing below is the complete code for the patching experiment of \cref{fig:residpatch} on a single pair of
prompts. It prints a logit difference of $3.36$ for the clean prompt and $-2.48$ for the corrupted one, and
then the effect at S2 and at the last position for each layer. The script \texttt{code/gpt2\_examples.py}
extends it to all figures of this tutorial.

\lstinputlisting[style=tutpython]{code/minimal_patching.py}

\begin{caution}
The listing was tested with TransformerLens 4.4. Older tutorials use the class \texttt{HookedTransformer},
which was removed in version 4.0 in September 2026. They run unchanged with an older version of the
library (\texttt{pip install "transformer-lens<4"}). The names of the hooks are the same in both.
\end{caution}

\subsection{Five first projects}

Reading about interpretability teaches less than doing it. The following projects are ordered by difficulty,
and each takes between an afternoon and a week.

\begin{enumerate}
\item \textbf{Break the running example.} Run the accompanying script. Then change the template: other verbs,
three names, the order of the clauses. Find out which of the figures of this tutorial change and which stay.
This is failure mode 8 of \cref{sec:pitfalls} in practice.
\item \textbf{Find the induction heads.} Feed GPT-2 small a sequence of random tokens followed by the same
sequence again. For every head, average the attention weight from each token of the second half to the token
after its first occurrence. A few heads will stand out. Compare them with \cref{tab:heads}.
\item \textbf{Solve the toy model.} Train the model \eqref{eq:toy} with $N=2$, $d=1$ for several values of
$\delta$ and compare the solution and its loss with \cref{prop:toy}. Then take $N=5$, $d=2$ and look for the
pentagon of \cref{fig:superposition}.
\item \textbf{Test a latent.} Pick a latent of a published dictionary on Neuronpedia. Write down a prediction
of where it fires and what it does when it is amplified. Then test both, including on text where it should
\emph{not} fire.
\item \textbf{Steer with controls.} Compute a difference of means \eqref{eq:diffmeans} for a simple contrast,
for example positive against negative reviews. Add it to the stream and measure the effect. Repeat with a random
direction of the same norm and with the direction found at another layer.
\end{enumerate}

\subsection{Open problems}

The field is young, and its foundations are still moving. Four questions are widely seen as open. The surveys
of \citet{sharkey2025open} and \citet{barez2026understanding} treat them at length.

\begin{itemize}
\item \textbf{What are the right units?} SAE latents are not canonical (\cref{sec:sae}). One alternative
decomposes the \emph{weights} of a network into simple parts instead of its activations
\citep{braun2025interpretability,bushnaq2025stochastic}. Another trains networks whose weights are sparse from
the start, which makes their circuits small and readable at some cost in capability \citep{gao2025weight}.
\item \textbf{How can an explanation be validated?} The tests of \cref{sec:circuits} work on narrow tasks. There
is no agreed standard for explanations of open-ended behaviour.
\item \textbf{What is the theory?} Why linear representations emerge, how much computation fits into
superposition, and when sparse dictionaries recover the true features are mathematical questions with partial
answers at best \citep{park2024linear,hanni2024mathematical,adler2024complexity}.
\item \textbf{What is it good for?} There is an active debate between two aims: the complete reverse engineering
of a network, and methods that are judged by their measurable use in tasks such as monitoring and auditing
\citep{nanda2025pragmatic,marks2025auditing}.
\end{itemize}

\subsection{What to read next}

\begin{itemize}
\item \citet{elhage2021mathematical} is the source of the residual-stream view of \cref{sec:transformer} and of
\cref{sec:pathexp}. Read it first.
\item \citet{elhage2022toy} develops superposition (\cref{sec:features}) with many experiments.
\item \citet{wang2023interpretability} is the model of a circuit analysis, and \citet{nanda2023progress} of a
complete one.
\item \citet{heimersheim2024use} is a practical guide to activation patching and its interpretation.
\item \citet{bricken2023monosemanticity} and \citet{templeton2024scaling} introduce sparse autoencoders, and
\citet{ameisen2025circuit} and \citet{lindsey2025biology} the attribution graphs built on them.
\item \citet{ferrando2024primer} and \citet{bereska2024mechanistic} are broad surveys. The glossary of
\citet{nanda2022glossary} explains the jargon of the field, and the ARENA course \citep{mcdougall2023arena}
teaches the methods through programming exercises.
\end{itemize}
